\documentclass{article}

% ready for submission
% \usepackage{neurips_2024}
\usepackage[preprint]{neurips_2024}

\usepackage[utf8]{inputenc}
\usepackage[T1]{fontenc}
\usepackage{hyperref}
\usepackage{url}
\usepackage{booktabs}
\usepackage{amsfonts}
\usepackage{nicefrac}
\usepackage{microtype}
\usepackage{xcolor}
\usepackage{graphicx}
\usepackage{amsmath}
\usepackage{multirow}
\usepackage{enumitem}

\title{Comparing Computer Vision Architectures for Pedestrian Counting at the Toomer's Corner Livestream}

\author{%
  Logan Bolton \\
  Auburn University \\
  \texttt{ldb0046@auburn.edu} \\
  \And
  Jase Schwanke \\
  Auburn University \\
  \texttt{jms0285@auburn.edu} \\
  \And
  Sathvik Prahadeeswaran \\
  Auburn University \\
  \texttt{srp0061@auburn.edu} \\
  \And
  Trey Wendell \\
  Auburn University \\
  \texttt{tdw0042@auburn.edu} \\
}

\begin{document}
\maketitle

\begin{abstract}
We study pedestrian counting from a single fixed outdoor camera — the
Toomer's Corner public livestream at Auburn University. Because no
labelled dataset exists for this scene, we build one: frames are
sampled from the livestream, pre-annotated with the open-vocabulary
SAM\,3 detector using the text prompt ``person,'' and then corrected
by human reviewers in Label Studio. We treat counting as an object
detection problem: a YOLO family detector outputs person boxes, and
the count is the number of boxes retained after a confidence
threshold. A pretrained YOLO26n serves as the zero-shot baseline, and
we fine-tune YOLO26s and YOLO26m variants on the curated data. A
large random hyperparameter search (195 successful trials across
eight sweep rounds) selects the final model by validation count MSE.
The fine-tuned YOLO26m reaches a per-image count MAE of 1.27 and MSE
of 2.79 on the held-out validation set, compared to a count MAE of
3.10 and MSE of 15.01 for the pretrained baseline --- roughly a
5\,$\times$ reduction in squared error from fine-tuning alone. The
main takeaway is that a small amount of in-domain labelled data,
combined with a modern open-vocabulary model for bootstrapping
labels, is enough to turn a generic pedestrian detector into a
reliable crowd counter for a fixed outdoor camera.
\end{abstract}

\section{Introduction}
Toomer's Corner is a public intersection at Auburn University that is
continuously streamed online. Traffic at the corner changes
dramatically by time of day and day of week --- from a handful of
pedestrians during class hours to dense crowds after sporting events
--- and an automatic crowd-counting system would let the university
quantify foot traffic without storing raw video. The livestream is,
however, shot from a single fixed rooftop angle with a wide
field-of-view, which is quite different from the clean frontal
pedestrian imagery that generic detectors are typically pretrained
on. This creates a mismatch: off-the-shelf detectors run on these
frames produce many missed detections, especially on distant or
partially-occluded pedestrians.

We therefore treat this as a concrete, scoped machine learning
problem: given a frame from the Toomer's Corner livestream, predict
the number of people visible in it. We compare a generic pretrained
detector to one fine-tuned on in-domain frames, measure how each
performs along accuracy and compute-efficiency axes, and use a large
random hyperparameter search to ensure the fine-tuned model is
actually well-tuned rather than accidentally configured.

\subsection{Problem Definition}
\begin{itemize}[leftmargin=1.5em]
    \item \textbf{Input:} a single RGB frame sampled from the
    Toomer's Corner livestream (typically $1280\times720$).
    \item \textbf{Output:} a non-negative integer $\hat{y}$
    estimating the number of people visible in the frame.
    \item \textbf{Why machine learning:} no closed-form feature can
    count pedestrians under real-world appearance variation
    (clothing, pose, lighting, occlusion, scale), so a learned
    detector is required.
    \item \textbf{Why challenging:} the fixed camera is far from
    many subjects, producing small-scale targets; crowds frequently
    self-occlude; lighting changes across the day (direct sun,
    shadow, night); and the scene contains easily-confused
    non-person objects such as mannequins in storefronts and
    reflections in shop windows.
\end{itemize}

\subsection{Project Scope and Objectives}
We intentionally constrained the project to a tractable scope: one
scene, one task (counting), one detector family (YOLO), with a
well-defined baseline vs.\ fine-tuned comparison. Specifically:
\begin{itemize}[leftmargin=1.5em]
    \item \textbf{Objective 1 --- Data:} build a $\geq$400-image
    labelled dataset of Toomer's Corner frames with person bounding
    boxes, using SAM\,3 to pre-annotate and Label Studio for human
    review.
    \item \textbf{Objective 2 --- Baseline:} evaluate an
    off-the-shelf pretrained YOLO26n on the same validation split
    as a zero-effort lower bound.
    \item \textbf{Objective 3 --- Fine-tuning:} fine-tune YOLO26s
    and YOLO26m on the curated data, and tune hyperparameters with
    a large random search whose primary metric is per-image count
    MSE (not standard mAP, because the downstream task is counting).
    \item \textbf{Objective 4 --- Analysis:} compare the two models
    on accuracy (count MSE, MAE), detection quality (mAP\textsubscript{50}),
    and compute (training time, inference latency), and surface
    qualitative failure modes.
\end{itemize}

\section{Dataset}
\subsection{Data Source}
Frames were captured from the public Toomer's Corner livestream by
\texttt{extract\_frames.py}, which runs \texttt{ffmpeg} against the
HLS URL at a fixed sampling rate (1~frame per second) and records
the wall-clock capture time in each filename. Sampling sessions
covered multiple days and times of day to diversify lighting and
crowd density. The resulting raw image directories
(\texttt{frames\_output/stream*/}) each contain several thousand
unlabelled frames; we subsampled roughly 500 frames for annotation,
discarding obvious duplicates (consecutive frames with no visible
motion) and frames with corrupted HLS segments.

\subsection{Annotation}
Manually drawing boxes for $\sim$15 people per image across 500
images is expensive, so we adopted a semi-automatic pipeline:
\begin{enumerate}[leftmargin=1.5em]
    \item \textbf{SAM\,3 pre-annotation.}
    \texttt{sam3/annotate.py} runs SAM\,3 (Segment Anything~3,
    an open-vocabulary detector) with the text prompt
    \texttt{"person"} and a confidence threshold of $0.25$.
    It writes YOLO-format label files (normalized
    \texttt{x\_center y\_center w h}) to \texttt{data/labels/raw/}.
    \item \textbf{Human review.}
    \texttt{sam3/review.py} launches a local Label Studio server,
    ingests every image with the SAM\,3 boxes as pre-loaded
    predictions, and four team members jointly correct them:
    adding missed detections, removing false positives (commonly
    mannequins, reflections, and tree trunks), and tightening
    boxes on distant pedestrians. Corrected labels are exported
    back in YOLO format.
    \item \textbf{Class schema.} A single class, \texttt{person},
    is used. Cyclists and riders are labelled as persons. Seated
    persons are labelled. Strollers, dogs, and vehicles are not.
\end{enumerate}
The pre-annotation step turned out to save a large fraction of
annotation time on sparse frames, but on dense frames reviewers
still had to add many missed detections by hand --- quantitatively,
SAM\,3's raw labels had a count MAE of 5.0 against the corrected
ground truth on the validation split, i.e.\ SAM\,3 alone is not
accurate enough for the downstream task.

\subsection{Dataset Statistics}
The final curated dataset (\texttt{data/manual\_label\_4\_16\_26\_v2})
is split 80/20 into train/validation at the image level by
\texttt{prepare\_dataset.py} with \texttt{seed=42}. Splits are
image-level (no frame from the same 1\,s neighbourhood is allowed
to straddle train and val, since consecutive frames are
near-duplicates). Basic statistics are summarised in
Table~\ref{tab:dataset_summary}.

\begin{table}[h]
    \centering
    \caption{Dataset summary. ``Avg.\ count'' is the mean number of
    person boxes per image. ``Max count'' is the single most
    crowded frame.}
    \label{tab:dataset_summary}
    \begin{tabular}{lcccc}
        \toprule
        Split & Images & Avg. count & Max count & Notes \\
        \midrule
        Train & 404 & 15.8 & 38 & Used for fine-tuning only \\
        Validation & 101 & 16.4 & 38 & Used for all reported metrics \\
        \bottomrule
    \end{tabular}
\end{table}

Per-image counts span 5 to 38, with most frames in the 10--20
range; true total across the validation split is 1659 people. The
data contains both sparse frames (individual pedestrians, clear
sidewalk) and dense frames (game-day crowds, multiple clusters).

\subsection{Data Quality and Preprocessing}
\textbf{Resolution.} Frames are stored at native livestream
resolution and passed to YOLO at \texttt{imgsz=1280} during both
training and inference. We briefly compared \texttt{imgsz=640} and
\texttt{960}; both degraded recall on distant/small pedestrians, so
\texttt{imgsz=1280} was fixed for the remainder of the sweep (see
Section~\ref{sec:sweep}).

\textbf{Filtering bad frames.} Frames with HLS corruption (gray
blocks, tear artefacts) were removed during annotation by the
reviewer.

\textbf{Leakage prevention.} Because consecutive 1\,FPS frames are
highly correlated, we avoided placing frames sampled within the
same short window on opposite sides of the train/val split. The
split is applied after filtering duplicates.

\textbf{Normalization and augmentation.} Pixel normalization is
delegated to Ultralytics' YOLO pipeline (scale to $[0,1]$, standard
BGR$\to$RGB). Training-time augmentation was treated as a
hyperparameter: mosaic, mixup, colour jitter (HSV), random scale,
and random erasing were all included in the search space, with
horizontal flip disabled (\texttt{fliplr=0}) because the scene has
asymmetric storefronts that a flipped model tends to hallucinate
boxes near. Close-mosaic was toggled on for the final epochs to
stabilize the final weights.

\section{Methods}
\subsection{Task Formulation}
We formulate counting as \emph{detection-then-count}: the model
produces a list of bounding boxes with confidence scores for the
\texttt{person} class, a confidence threshold~$\tau$ is applied,
and the count is simply the number of retained boxes. At
evaluation time we sweep $\tau \in
\{0.10, 0.15, 0.20, 0.25, 0.30, 0.35, 0.40, 0.50, 0.60\}$
and report the threshold that minimises count MSE on the
validation split, i.e.\ the threshold is part of the evaluation,
not the model. We chose detection-then-count over direct
regression because it yields interpretable box-level outputs that
the reviewers could inspect to debug failure modes, and because
the SAM\,3 labelling pipeline produces boxes natively.

\subsection{Baseline Model: Pretrained YOLO26n}
The baseline is the smallest (nano) YOLO26 checkpoint,
\texttt{yolo26n.pt}, loaded with no fine-tuning and evaluated on
the validation split with \texttt{classes=[0]} (COCO person class).
It represents what one would get with zero labelling effort and is
the natural ``is fine-tuning worth the cost'' reference point.
Nothing is trained; only the inference-time confidence threshold
is tuned on the validation set.

\subsection{Fine-tuned Model: YOLO26m}
The main model is \texttt{yolo26m.pt} (the medium YOLO26
checkpoint) fine-tuned on our curated training split via
\texttt{ultralytics.YOLO.train()}. All layers are unfrozen; we
start from COCO-pretrained weights. Person class only (class
index~0). We also trained YOLO26s variants as a middle datapoint;
YOLO26n fine-tuning was dropped early in the sweep because every
\texttt{yolo26n} trial in sweep~v5 produced catastrophic count
errors, suggesting the nano backbone is too small to resolve
distant pedestrians at this scene's scale.

\subsection{Third Model Explored: Moondream 2 (VLM)}
We also evaluated the Moondream~2 vision-language model
(\texttt{vikhyatk/moondream2}, revision \texttt{2025-06-21}) via
its \texttt{detect(image, "person")} interface. This serves as an
``easy knob'' reference: no training, no dataset, open-vocabulary
prompt. In our qualitative runs on representative Toomer's Corner
frames, Moondream~2 produced far fewer detections than the true
crowd count (it tended to localise only the most salient
foreground pedestrians), so it is not competitive with either
YOLO setup for dense-scene counting; we report it only as
qualitative context.

\subsection{Training Setup}
Fine-tuning is driven by \texttt{train.py} (for single runs) and
\texttt{sweep.py} (for hyperparameter search). Per-trial defaults
during the search: 50 epochs, early stopping patience 15, image
size $1280$, batch 16, Ultralytics dataloader with
\texttt{cache=ram}. Hardware: one NVIDIA GPU per worker process,
run with two workers in parallel
(\texttt{CUDA\_VISIBLE\_DEVICES=0,1}) so that two trials train
concurrently without contending over the same device. Model
weights, optimizer type, learning rate schedule, loss weights,
and all augmentation probabilities were \emph{not} hand-picked ---
they are hyperparameters (Section~\ref{sec:sweep}).

Best trial hyperparameters (sweep~v10, trial
\texttt{1776384144137\_w0\_0015}):
\texttt{model=yolo26m.pt},
\texttt{imgsz=1280}, \texttt{batch=16}, \texttt{optimizer=auto},
\texttt{lr0=0.01}, \texttt{lrf=0.1}, \texttt{momentum=0.95},
\texttt{weight\_decay=1e-3}, \texttt{warmup\_epochs=1},
\texttt{cos\_lr=True}, \texttt{box=10.0}, \texttt{cls=1.0},
\texttt{dfl=1.5}, HSV~\texttt{(0.0, 0.5, 0.4)},
\texttt{scale=0.25}, \texttt{fliplr=0.0}, \texttt{flipud=0.1},
\texttt{mosaic=0.5}, \texttt{mixup=0.25},
\texttt{close\_mosaic=5}, \texttt{dropout=0.2}.

\subsection{Evaluation Metrics}
\label{sec:metrics}
Because counting is the end task, our primary metrics are
per-image count errors, not box-level localisation:
\begin{itemize}[leftmargin=1.5em]
    \item \textbf{Count MSE} (primary sweep metric):
    $\tfrac{1}{N}\sum_i (\hat{y}_i - y_i)^2$, where $\hat{y}_i$ is
    the number of surviving boxes and $y_i$ the ground-truth count.
    \item \textbf{Count MAE:}
    $\tfrac{1}{N}\sum_i |\hat{y}_i - y_i|$, reported because
    absolute-error units are easy to reason about
    (``off by $k$ people per frame'').
    \item \textbf{Count bias:} mean signed error, which tells us
    whether a model systematically under- or over-counts.
    \item \textbf{mAP\textsubscript{50}:} standard COCO-style mean
    average precision at IoU\,=\,0.5, reported for context
    (box-quality, not counting-quality).
    \item \textbf{Training time} per trial (seconds), recorded
    automatically by the sweep.
    \item \textbf{Inference latency} (seconds) for the full
    validation split on a single GPU, end-to-end including image
    decode.
\end{itemize}

\section{Experiments}
\subsection{Main Comparison}
The baseline (pretrained YOLO26n) and fine-tuned YOLO26m are
evaluated on the same 101-image validation split using the same
preprocessing. For each model, the confidence threshold is tuned
on the validation set by grid-search to minimise count MSE, then
the per-image count errors at that threshold are reported in
Section~\ref{sec:quant}.

\subsection{Hyperparameter Sweep}
\label{sec:sweep}
\texttt{sweep.py} implements a random search over a 26-dimensional
space (model size, optimizer, LR schedule, loss weights,
augmentation probabilities, regularization). Each trial samples
one config, trains one YOLO model to completion with that config,
evaluates the best-confidence-threshold count MSE on the
validation split, and appends a row to
\texttt{sweeps/<name>/results.csv}. We ran eight sweep rounds
(\texttt{v1, v3, v4, v5, v7, v8, v9, v10}) for a total of 236
trials, 195 of which completed successfully --- the remainder died
on out-of-memory or transient CUDA errors. Between rounds, we
pruned the search space using the analysis of per-axis mean count
MSE from the previous round (e.g.\ after v5, we dropped
\texttt{yolo26n} because every nano trial regressed badly, and
pinned several augmentation knobs whose variance across the sweep
was dominated by noise). The final sweep \texttt{v10} is therefore
a tighter, focused search over the parameters that still showed
meaningful spread.

\subsection{Trend Analysis}
Frame filenames record the wall-clock capture time (see
\texttt{extract\_frames.py}), so applying the final fine-tuned
model to unlabelled frames from a multi-hour stream produces a
per-second count time series. We report one such time series in
Section~\ref{sec:trend}.

\section{Results}
\subsection{Quantitative Results}
\label{sec:quant}
Table~\ref{tab:main_results} reports the head-to-head comparison
on the validation split at each model's best confidence threshold.
Fine-tuning reduces count MSE by roughly a factor of five and
halves count MAE. The fine-tuned model is also nearly unbiased
($-0.07$ people/frame), while the baseline systematically
under-counts ($-1.63$).

\begin{table}[t]
    \centering
    \caption{Main validation-set results. Arrows indicate direction
    of improvement. Training and inference times are recorded on the
    same single-GPU setup. The baseline is \emph{not} trained, so
    its training time is zero.}
    \label{tab:main_results}
    \begin{tabular}{lccccc}
        \toprule
        Model & MAE $\downarrow$ & MSE $\downarrow$ & mAP\textsubscript{50} $\uparrow$ & Train time & Inference (full val) \\
        \midrule
        Pretrained YOLO26n (baseline) & 3.10 & 15.01 & --- & 0\,s & 1.3\,s / 91 imgs \\
        Fine-tuned YOLO26m (ours)     & \textbf{1.27} & \textbf{2.79} & \textbf{0.908} & 588\,s & 2.9\,s / 91 imgs \\
        \bottomrule
    \end{tabular}
\end{table}

Hyperparameter sweep summary (all 195 successful trials): Count
MSE ranged from 2.79 (best) to 20.8 (worst); every
\texttt{yolo26m} fine-tuning run in the final round was better
than the pretrained baseline. The best-per-round MSE (v1
$\to$ v10) dropped from 7.06 to 2.79, indicating the search-space
pruning strategy worked: later rounds reliably produced better
configurations than earlier ones.

\subsection{Compute Efficiency Results}
Fine-tuning costs a one-time $\sim$10 minutes on a single GPU at
our training settings. Inference cost is $\sim$30\,ms/frame for
the fine-tuned medium model vs.\ $\sim$14\,ms/frame for the nano
baseline --- in absolute terms, both are more than fast enough to
run in real time on a 1\,FPS livestream. The full $26$-dim sweep
across 195 trials consumed roughly $5$--$7$ minutes per trial on
average, parallelised across two GPUs, for a total wall-clock cost
under one day of compute.

\begin{table}[t]
    \centering
    \caption{Compute efficiency comparison. Training time is
    per-trial; inference is per-frame at \texttt{imgsz=1280} on a
    single GPU.}
    \label{tab:compute_efficiency}
    \begin{tabular}{lcc}
        \toprule
        Model & Train time (best trial) & Inference latency \\
        \midrule
        Pretrained YOLO26n (baseline) & --- & $\sim$14\,ms/frame \\
        Fine-tuned YOLO26m (ours)     & $\sim$10\,min & $\sim$32\,ms/frame \\
        \bottomrule
    \end{tabular}
\end{table}

\subsection{Trend Analysis Results}
\label{sec:trend}
Applying the fine-tuned model to an unlabelled 1\,FPS capture of
Toomer's Corner recovers an intuitive crowd-count time series:
counts are low during mid-morning class hours, rise steadily at
lunchtime, and peak during the late afternoon as classes let out.
This is not a rigorous longitudinal study --- a multi-day capture
would be required to separate time-of-day effects from weekday
vs.\ weekend effects --- but it confirms that the predicted counts
behave as a crowd-density estimator should.

\subsection{Qualitative Examples}
Figure~\ref{fig:qualitative} shows two representative cases. In
sparse scenes, both models detect the visible pedestrians
correctly; the gap between baseline and fine-tuned is invisible
to the eye. In dense scenes, the baseline misses many distant or
partially-occluded pedestrians, while the fine-tuned model
recovers most of them. Failure modes for the fine-tuned model
cluster around three cases: (1) extremely small pedestrians near
the camera vanishing point, (2) heavy self-occlusion in tightly
packed groups, where one bounding box often covers two people,
and (3) nighttime frames, where both false positives (reflections,
signage) and false negatives (underlit pedestrians) increase.

\begin{figure}[t]
    \centering
    \fbox{\rule{0pt}{2.2in}\rule{0.95\linewidth}{0pt}}
    \caption{Qualitative predictions. Left: sparse daytime frame,
    both models correct. Right: dense crowd, baseline misses many
    distant pedestrians while the fine-tuned model recovers them.}
    \label{fig:qualitative}
\end{figure}

\section{Discussion}
\subsection{What Worked Well}
The single strongest finding is that a relatively small curated
dataset ($\sim$500 frames) is enough to turn a generic pedestrian
detector into a reliable crowd counter for this scene. Two design
choices mattered: (i)~SAM\,3 pre-annotation cut the labelling
effort enough that four students could realistically produce the
dataset in a class-project time budget; and (ii)~treating the
confidence threshold and all augmentation knobs as search
hyperparameters, rather than hand-tuning them, reliably surfaced
configurations noticeably better than any default we tried by
hand. Making \emph{count MSE} (not mAP) the sweep's selection
metric also mattered --- we observed several trials with higher
mAP\textsubscript{50} but worse counting, because mAP rewards
tight localisation but not box-count fidelity.

\subsection{Limitations}
\begin{itemize}[leftmargin=1.5em]
    \item \textbf{Single scene.} The model is fine-tuned on one
    fixed camera angle. We do not expect it to transfer to other
    streams without re-labelling.
    \item \textbf{Label noise.} Because annotation is based on
    SAM\,3 pre-labels, systematically hard cases for SAM\,3
    (e.g.\ very distant or low-contrast pedestrians) are
    under-represented in the human-added corrections, and some
    residual labelling noise survives into the training set.
    \item \textbf{Small test set.} With only 101 validation
    images, count MSE differences below roughly $0.3$ are within
    trial-to-trial noise, so we cannot reliably distinguish the
    top few sweep trials from one another.
    \item \textbf{Day-dominated data.} Most captured frames are
    daytime; the model is noticeably less reliable at night
    (see below).
    \item \textbf{Counting via boxes.} The detection-then-count
    pipeline inherits all standard detection failure modes (NMS
    merging two people into one box, double-counting a person
    split across occlusion), which a density-map approach might
    sidestep --- we did not compare to such a method.
\end{itemize}

\subsection{Sources of Error and Bias}
Partial people at the frame border are sometimes labelled and
sometimes not, which reviewers tried to handle consistently but
could not fully eliminate; this is a measurable source of
labeller disagreement. Dense crowds produce near-duplicate
bounding boxes post-NMS that can collapse multiple nearby people
into a single detection. Weather and night-time frames were
deliberately kept in the set but are under-represented, meaning
the model's night-time count variance is higher than daytime.
Finally, because the SAM\,3 pre-annotator biases what reviewers
see first, there is a small but nonzero risk of confirmation bias
in the labels --- reviewers are more likely to approve an
existing box than to notice a genuinely missed pedestrian in a
visually busy region.

\section{Conclusion}
We built a person-counting system for the Toomer's Corner
livestream end-to-end: frame extraction, semi-automatic
annotation with SAM\,3 + Label Studio, fine-tuning in the YOLO26
family, and a large hyperparameter search whose primary metric is
count MSE rather than mAP. Fine-tuning cut per-image count MAE
from 3.10 (pretrained YOLO26n) to 1.27 (fine-tuned YOLO26m), and
count MSE from 15.01 to 2.79, on the same held-out validation
split. The fine-tuned model is essentially unbiased; the baseline
systematically under-counts. Compute costs are negligible
(one-time $\sim$10\,min fine-tuning, $\sim$30\,ms/frame
inference), so this approach is viable as an always-on counter.
Natural next steps include multi-scene transfer, density-map
methods for very dense frames, and a larger longitudinal capture
to quantify time-of-day and weekday-vs-weekend effects rigorously.

\section{Team Contributions}
\begin{itemize}[leftmargin=1.5em]
    \item \textbf{Jase Schwanke:} Frame extraction pipeline
    (\texttt{extract\_frames.py}), livestream sampling strategy,
    qualitative failure-case curation, writing contributions to
    the dataset and discussion sections.
    \item \textbf{Logan Bolton:} SAM\,3 annotation pipeline
    (\texttt{sam3/annotate.py}, \texttt{sam3/review.py}) and
    Label Studio integration, \texttt{sweep.py} implementation
    and multi-GPU orchestration, final sweep analysis.
    \item \textbf{Sathvik Prahadeeswaran:} Baseline YOLO
    evaluation, Moondream~2 VLM exploration, hyperparameter
    search-space pruning between sweep rounds, compute-efficiency
    measurements.
    \item \textbf{Trey Wendell:} Dataset splitting
    (\texttt{prepare\_dataset.py}), annotation review shifts,
    \texttt{compare\_models.ipynb} for final cross-model plots,
    writing contributions to the introduction, methods, and
    conclusion.
\end{itemize}
All authors participated in labelling shifts in Label Studio and
contributed to reviewing the final report.

\section*{Code and Reproducibility}
The project repository contains everything needed to reproduce
the results. Entry points:
\begin{itemize}[leftmargin=1.5em]
    \item \texttt{extract\_frames.py} --- sample frames from the
    livestream.
    \item \texttt{sam3/annotate.py} --- generate raw SAM\,3 labels.
    \item \texttt{sam3/review.py} --- launch Label Studio for
    human correction.
    \item \texttt{prepare\_dataset.py} --- build train/val split
    and \texttt{dataset.yaml}.
    \item \texttt{train.py} --- single fine-tuning run.
    \item \texttt{sweep.py} --- multi-GPU random hyperparameter
    search. Results land in \texttt{sweeps/<name>/results.csv}.
    \item \texttt{compare\_models.ipynb} --- side-by-side count
    metrics, best-confidence-threshold selection, and plots used
    in this report.
\end{itemize}
Environment is managed with \texttt{uv}
(\texttt{pyproject.toml}~+~\texttt{uv.lock}); a single
\texttt{uv sync} pins all dependencies. All randomness uses fixed
seeds (\texttt{prepare\_dataset.py} seed~$=42$; sweep workers
seed from a reproducible base). Best model weights for each
retained trial are stored under
\texttt{sweeps/<name>/trials/<trial\_id>/best.pt}. See the
repository \texttt{README.md} for a quickstart.

\section*{References}
\begin{itemize}[leftmargin=1.5em]
    \item Ultralytics, \emph{YOLO26}. \url{https://docs.ultralytics.com/}
    \item A. Kirillov et al., \emph{Segment Anything}, ICCV 2023; SAM\,3 open-vocabulary extension.
    \item Label Studio (Heartex), \url{https://labelstud.io/}
    \item Moondream team, \emph{Moondream 2}. \url{https://huggingface.co/vikhyatk/moondream2}
    \item T.-Y. Lin et al., \emph{Microsoft COCO: Common Objects in Context}, ECCV 2014.
\end{itemize}

\end{document}
